% Generated draft: requires mathematical and source review.
\documentclass[11pt]{article}
\usepackage[margin=25mm]{geometry}
\usepackage{fontspec}
\setmainfont{Latin Modern Roman}
\usepackage{amsmath,amssymb,mathtools,unicode-math}
\setmathfont{Latin Modern Math}
\usepackage{enumitem,needspace,xurl,hyperref}
\setlength{\parindent}{0pt}
\setlength{\parskip}{5pt}
\setlength{\emergencystretch}{3em}
\newcommand{\N}{\mathbb{N}}
\newcommand{\Z}{\mathbb{Z}}
\newcommand{\Q}{\mathbb{Q}}
\newcommand{\R}{\mathbb{R}}
\newcommand{\C}{\mathbb{C}}
\newcommand{\F}{\mathbb{F}}
\newcommand{\A}{\mathbb{A}}
\newcommand{\PP}{\mathbb{P}}
\newcommand{\E}{\mathbb{E}}
\newcommand{\eps}{\varepsilon}
\newcommand{\dd}{\,\mathrm{d}}
\begin{document}
\section*{UnsolvedMath ID: 701493}
\textbf{Source problem code:} catalogue entry
\subsection{Definitions and mathematical statement}
Let \(G\) be a finite group and let \(p\) be a prime. Write \(\operatorname{Irr}(G)\) for the set of irreducible complex characters of \(G\). Let \(B\) be a \(p\)-block of \(G\), in the usual sense of modular representation theory, and write \(\operatorname{Irr}(B)\) for the ordinary irreducible characters belonging to \(B\). For \(\chi\in\operatorname{Irr}(G)\), its degree is \(\chi(1)\). An element \(g\in G\) is \(p\)-singular if \(p\) divides its order.

The conjecture attributed in the supplied record to K. Harada asserts that, for every such \(G,p,B\) and every nonempty subset \(J\subseteq\operatorname{Irr}(B)\),
\[
\left(\forall g\in G\text{ with }p\mid\operatorname{ord}(g),\quad
\sum_{\chi\in J}\chi(1)\chi(g)=0\right)
\quad\Longrightarrow\quad J=\operatorname{Irr}(B).
\]
\subsection{Short English statement}
Harada's conjecture states that a nonempty subset of the irreducible complex characters in one \(p\)-block must be the entire block if its degree-weighted character sum vanishes on every \(p\)-singular element.
\subsection{Sources}
\begin{itemize}
\item Local UnsolvedMath registry, permanent ID \(701493\), exact source problem code \texttt{catalogue entry}, title: Kourovka Notebook Problem 18.64. Supplied upstream revision: \texttt{d371ee917ac9647fbb06cd6fc529deec6f3dda82}.
\item \emph{The Kourovka Notebook (Unsolved Problems in Group Theory)}, Problem 18.64. The supplied background identifies the 18th issue (2014) and the 21st issue edition (2026), with identifier \texttt{arXiv:1401.0300v48}. The collection description names E. I. Khukhro and V. D. Mazurov as editors.
\item The registry lists V. D. Mazurov as proposer and 2014 as the proposed year; the statement itself attributes the conjecture to K. Harada.
\item The supplied registry labels the problem open and records an update date of 2026-10-16. This is supplied source metadata, not an independently verified assessment of its mathematical status.
\end{itemize}
\end{document}
